\chapter{Conclusion: What the Half Explains and What It Does Not}
\label{chap:conclusion}

The number $1/2$ is not an arbitrary decoration of the zeta function. It is the exact centre of the complement $s\mapsto1-s$ and the fixed line of the anti-linear involution $s\mapsto1-\bar s$. In a normalized binary model, it is also the unique point of maximum Shannon entropy, the unique equal-magnitude condition for two-channel destructive interference, the Bloch equator, and the unique phase fraction producing the sign $-1$ under $e^{2\pi i\sigma}$.

These facts are exact. They can be summarized by the binary-spinor equivalence theorem:
\begin{align*}
 \text{self-complementarity}
 &\iff \text{maximum binary entropy}\\
 &\iff \text{spinorial sign}\\
 &\iff \text{phase-locked cancellation}\\
 &\iff \sigma=\frac12.
\end{align*}
The Dirichlet eta function then supplies a classical analytic bridge:
\[
 \eta(1)=\ln2,
 \qquad
 \eta(s)=(1-2^{1-s})\zeta(s),
\]
and eta and zeta have the same zeros in the open critical strip. Binary alternation is therefore genuinely present in the analytic object whose zeros are at issue.

What has not been proved is the implication from an eta or xi zero to the normalized two-channel cancellation event. The infinite alternating series can cancel without reducing to two equal-norm components. The odd and even analytic channels both vanish at a zeta zero and therefore do not define a normalized state there. Functional symmetry can pair off-axis zeros without fixing them. Entropy selects a line but not the zero ordinates. Spinorial phase supplies the correct sign but not the canonical zeta loop.

The unresolved bridge can be stated cleanly: construct an analytic, involution-covariant, non-vanishing Hilbert-space vector whose fixed detector projection equals $\xi(s)$ up to a non-zero factor, and whose complementary component norms are $\Re s$ and $1-\Re s$ in a canonical swap-invariant metric. If this is done, the cancellation lemma proves RH. If the construction is further derived from a positive kernel or self-adjoint operator, the information-spinor picture becomes part of a stronger spectral mechanism.

The monograph therefore reaches a definite but bounded conclusion:
\begin{quote}
The half-axis has a unified structural explanation inside a precise binary-spinor model, and the eta function makes the binary connection to zeta zeros exact. The Riemann Hypothesis is reduced, within this programme, to a canonical representation problem that remains open.
\end{quote}

This is not the end of the derivation but the point where derivation must replace analogy. The next progress will not come from repeating that $\ln2$, $\pi$, and $1/2$ appear together. It will come from constructing the kernel, metric, or operator that forces them to appear together at every non-trivial zero.

\status{Research monograph conclusion. RH is not claimed as proved. The exact results and the blocking bridge are explicitly separated.}
